\subsection{Leaf recurrence}
%
Let $i\in V(Q)$ be a leaf vertex of a quiver $Q$, connected to some other vertex $j\in V(Q)$. We call $i$ a \emph{\Louiseleaf} if both quivers $Q-\{i\}$ and $Q-\{i,j\}$ are \Louise.  This automatically implies that $Q$ itself is \Louise.

%

\begin{proposition}\label{prop:Louise_leaf}
Fix a field $\Ga$, and consider cluster varieties over $\Ga$.
Let $\Qice$ be a really full rank ice quiver with $m$ frozen vertices and mutable part $Q$. Let $i\in V(Q)$ be a \Louiseleaf in $Q$ connected to $j\in V(Q)$. Then the subvariety of $\XQice$ given by $x_i=0$ is isomorphic to $\Ga\times \Xcal(\Qice')$, where $\Qice'$ is a really full rank quiver with $m$ frozen vertices and mutable part $Q':=Q-\{i,j\}$.
\end{proposition}
\begin{proof}%
Our goal is to construct a quiver $\Qice'$ satisfying the above assumptions and an isomorphism
\begin{equation}\label{eq:leaf1}
  \{\bx\in\Xcal(\Qice)\mid x_i=0\} \cong \F\times \Xcal(Q').
\end{equation}

  By Proposition~\ref{prop:Mulcover}, the subvariety of $\Xcal(\Qice)$ given by $x_i=0$ is contained inside the subvariety of $\XQice$ given by $x_j\neq0$. The quiver $Q-\{j\}$ is the disjoint union of $Q'$ and a single isolated vertex. By assumption, $Q'$ is \Louise and thus so is $Q-\{j\}$. By Proposition~\ref{prop:Mulloc}, we have $\AQice[x_j^{-1}]\cong \AX(\Qicefr[j])$. We therefore find
\begin{equation}\label{eq:leaf2}
    \{\bx\in\Xcal(\Qice)\mid x_i=0\} \cong \{\bx\in\Xcal(\Qicefr[j])\mid x_i=0\}.
\end{equation}

%
Let $A= \AX(\Qicefr[j] - \{i\})[x_i^{\pm 1}]$ be obtained from the cluster algebra $ \AX(\Qicefr[j] - \{i\})$ by adjoining $x_i$ and its inverse.  The cluster algebra $\AX(\Qicefr[j])$ is isomorphic 
%
 to the subalgebra of $A$ generated by $\AX(\Qicefr[j] - \{i\})$, the cluster variable $x_i$ and the mutation $x'_i$, which is of the form $f/x_i$ where $f \in \AX(\Qicefr[j] - \{i\})$.  

Let $\Xcal':=\Xcal(\Qicefr[j] - \{i\})$.  It follows that 
\begin{equation}\label{eq:leaf2.5}
  \Xcal(\Qicefr[j])\cong \{\left((x_i,x_i'),\bx'\right)\in\Ga^2\times\Xcal'\mid x_ix_i'=M+x_jM'\},
\end{equation}
where $M,M'$ are monomials in the frozen variables of $\Qicefr[j] - \{i\}$. We therefore see that the right-hand side of~\eqref{eq:leaf2} is given by 
\begin{equation}\label{eq:leaf3}
  \{\bx\in\Xcal(\Qicefr[j])\mid x_i=0\}
\cong \{\left((x_i,x_i'),\bx'\right)\in\Ga^2\times\Xcal'\mid x_i=0 \text{ and } x_ix_i'=M+x_jM'\}.
\end{equation}
We see that the variable $x_i'$ is free, and thus the right-hand side of~\eqref{eq:leaf3} is isomorphic to the direct product of $\Ga$ and 
\begin{equation} \label{eq:yM'/M=1} 
\left\{\bx'\in \Xcal'\ \middle|\  \frac{x_jM'}{M}=-1\right\}.
\end{equation}
%
%
%
%
%
%
%
%
%
%
%
%
%
%
%
  It remains to show that the locus~\eqref{eq:yM'/M=1} is isomorphic to $\Xcal(\Qice')$ for some quiver $\Qice'$ satisfying the conditions in \cref{prop:Louise_leaf}. By assumption, $\Bice:=\BQice$ is really full rank. Thus, there exist integers $(\alpha_k)_{k\in[n+m]}$ such that 
  \begin{equation*}%
    \sum_{k} \alpha_k \bice_k=e_i,
  \end{equation*}
where $\bice_k$ denotes the $k$-th row of $\Bice$ and $e_i=(0,\dots,0,1,0,\dots,0)$ is the $i$-th standard basis vector in $\Z^n$. 

%
Since $i$ is a leaf, the row $\bice_i$ has a single nonzero entry in the $j$-th column. Let $\Bice\pA$ be obtained from $\Bice$ by removing the $i$-th row and the $j$-th column.\footnote{When removing rows and columns from matrices, we preserve the labels of the remaining rows and columns. For instance, removing row $1$ from $\Bice$ produces a matrix with rows labeled $2,\dots,n+m$.} Then we have
  \begin{equation*}%
    \sum_{k\neq i} \alpha_k \bice\pA_k=\bar e_i,
  \end{equation*}
where $\bar e_i$ is the $i$-th standard basis vector in $\Z^{n-1}$. Moreover, the matrix $\Bice\pA$ is still really full rank.
%
%
%
%
%

Let $\Vfro:=[n+1,n+m]$ be the set of frozen vertices of $\Qice$. Then the set of frozen vertices of $\Qicefr[j]-\{i\}$ is $\Vfro\sqcup\{j\}$.  We have $\gcd(\alpha_k\mid k\in \Vfro\sqcup\{j\})=1$
%
 since all the rows with a nonzero entry in column $i$ of $\Bice$ belong to $\Vfro\sqcup\{j\}$.  By \cite[Section 5]{LS}, replacing a frozen row by its negative, or adding a frozen row to another frozen row produces an isomorphic cluster variety.  After applying a series of such transformations, we obtain a really full rank $(n+m-1)\times(n-1)$ exchange matrix $\Bice\pB$ and a family of integers $(\alpha\pB_k)_{k\neq i}$ satisfying 
   \begin{equation}\label{eq:alpha_kBice''_k=e_i}
    \sum_{k\neq i} \alpha\pB_k \bice\pB_k=\bar e_i,
  \end{equation}
such that moreover for some $\ell\in\Vfro\sqcup\{j\}$ we have $\alpha\pB_\ell=1$. 

Now, let $\Bice\pC$ be the $(n+m-1)\times(n-2)$ matrix obtained from $\Bice\pB$ by further removing the $i$-th column. Thus, $\AX(\Bice\pC)\cong \AX(\Qicefr[j] - \{i\})$, and by construction, $\sum_{k\neq i} \alpha\pB_k\bice\pC_k=0$ is the zero vector in $\Z^{n-2}$.  Thus, by~\cite[Proposition 5.1]{LS}, the numbers $(\alpha\pB_k)_{k\neq i}$ describe a one-parameter subgroup $\Gm$ of the cluster automorphism torus $T=\Aut(\AX(\Bice\pC))$ (see \cite[Section 5]{LS}), with $z\in\Gm$ acting by $x_k\mapsto z^{\alpha\pB_k}x_k$. Equation~\eqref{eq:alpha_kBice''_k=e_i} implies that $z$ acts on the monomial $\frac{x_jM'}{M}$ by $z^{\pm 1}$. Thus, every $\Gm$-orbit contains a unique point in the locus~\eqref{eq:yM'/M=1}. In other words, the locus~\eqref{eq:yM'/M=1} is isomorphic to the quotient $\Xcal(\Bice\pC)/\Gm$.  Since we assumed that $\alpha\pB_\ell = 1$, the quotient $\Xcal(\Bice\pC)/\Gm$ can be obtained by setting $x_\ell=1$.  Let $\Bice'$ be the $(n+m-2)\times(n-2)$ matrix obtained from $\Bice\pC$ by removing the $\ell$-th row. Let $\Qice'$ be the associated quiver, with $m$ frozen vertices corresponding to the rows $(\Vfro\sqcup\{j\})\setminus\{\ell\}$ and mutable part $Q-\{i,j\}$. Clearly, $\Qice'$ is still really full rank. We find that indeed the locus~\eqref{eq:yM'/M=1} is isomorphic to $\Xcal(\Qice')$.
%
%
  %
  %
\end{proof}

\begin{remark}
\Cref{prop:Louise_leaf} generalizes to the case of $i$ being either a source or a sink in $Q$.
\end{remark}
The following result will be later compared to the \FLY skein relation~\eqref{eq:HOMFLY_dfn}.
\begin{corollary}\label{cor:RQ_skein}
  Let $Q$ be really full rank, and let $i\in V(Q)$ be a \Louiseleaf in $Q$ connected to $j\in V(Q)$. Let $Q':=Q-\{i\}$ and $Q'':=Q-\{i,j\}$. Then
  \begin{equation}\label{eq:RQ_skein}
    \RQ(q)=(q-1)\RQX{Q'}{q} + q\RQX{Q''}{q}.
  \end{equation}
\end{corollary}
\begin{proof}
  Let $\Qice$ be really full rank with mutable quiver $Q$, and $\Xcal:=\XQice$. Let $\Ucal:=\{x_i\neq0\}$ and $\Zcal:=\{x_i=0\}$ be two subvarieties of $\Xcal$. Then $\Ucal$ is a really full rank cluster variety of type $Q'$ and $\Zcal$ is described in \cref{prop:Louise_leaf}. The result follows since
  \begin{equation*}%
    \#\Xcal(\F_q)=\#\Ucal(\F_q) + \#\Zcal(\F_q).    \qedhere
  \end{equation*}
\end{proof}

\section{Quivers, links, and plabic graphs}
In this section, we study some basic relations between the properties of plabic graphs and the associated quivers and links.
 %

\subsection{Plabic graph quivers}\label{sec:plabquiver}
Recall the background on plabic graphs from \cref{sec:intro:plabic_link,sec:plabic-graphs}. We continue to assume that each plabic graph $G$ is trivalent, and that each interior face of $G$ is simply connected.
%
%
%
%
%
%
%
%
%
%
%
%

We say that two plabic graphs are \emph{move equivalent} if they are connected by a sequence of moves in \cref{fig:moves}. For the square move in \figref{fig:moves}(b), we require that the four faces $A,B,C,D$ adjacent to the square face in clockwise order satisfy $A\neq B\neq C\neq D\neq A$. (Having $A=C$ or $B=D$ is allowed.) Cf.~\cite[Restriction~6.3]{FPST}; this restriction is needed in order for natural statements such as \cref{prop:simple_vs_moves} to hold.

\begin{proposition}\label{prop:trivred} 
%
Let $G$ be a connected plabic graph. Then $G$ can be transformed using only tail removal (\figref{fig:moves}(c)) into:
\begin{itemize}
\item the graph in \figref{fig:trivred}(a) if $G$ has no interior faces;
\item the graph in \figref{fig:trivred}(b) or \figref{fig:trivred}(c) if $G$ has one interior face;
\item a plabic graph with no boundary vertices, if $G$ has two or more interior faces.
\end{itemize}
We call this graph the \emph{tail reduction} of $G$. %
\end{proposition}
\begin{proof}
Let $G$ be a connected plabic graph.  By repeatedly applying tail removal  to $G$ (without changing connectedness or the number of interior faces), we may assume no more tail removals are possible.  If $G$ has no boundary vertices then $G$ must have two or more interior faces.  If $G$ has a boundary vertex $b$ connected to an interior vertex $v$, then $v$ must be incident to a loop edge, for otherwise we may apply a tail removal to $b$.  In this case, $G$ has one interior face.  Finally, if $G$ has no interior vertices, then $G$ consists of a single edge connecting two boundary vertices, and has no interior faces.
\end{proof}

\begin{figure}
  \includegraphics[width=0.6\textwidth]{figures/trivred}
  \caption{\label{fig:trivred} Tail reductions that have boundary vertices.}
\end{figure}

\begin{proposition}\ \label{prop:simple_vs_moves}
\begin{itemize}
\item Suppose that two simple plabic graphs $G$ and $G'$ are move equivalent. Then $\QG$ and $\QX(G')$ are mutation equivalent.  
\item Any plabic graph $G'$ move equivalent to a simple plabic graph $G$ is also simple.
\end{itemize}
\end{proposition}
\begin{proof}
Direct check.
\end{proof}

%
%
%
%
%
%

%
%
%
%
%
%

\begin{figure}
  \includegraphics[width=0.65\textwidth]{figures/interior_leaves}
  \caption{\label{fig:interior_leaves} The behavior of the link $\Lplab_G$ at interior leaves.}
\end{figure}

Recall from \cref{sec:intro:quivers-point-count} that each plabic graph gives rise to a link $\Lplab_G$ and that our main conjecture (\cref{conj:main}) yields a relation between the polynomials $\RQG(q)$ and $P(\Lplab_G;a,z)$.  For a rational function $R(q) = P(q)/Q(q)$, the \emph{leading coefficient} of $R(q)$ is defined as the ratio of the coefficient of $q^{\deg P}$ in $P(q)$ and the coefficient of $q^{\deg Q}$ in $Q(q)$. While \cref{conj:main} applies only to simple plabic graphs (cf. \cref{sec:counterex}), the following more basic statement appears to hold for arbitrary plabic graphs.  In particular, we allow plabic graphs with interior leaves, though we still insist that interior faces are simply connected.  The link $\Lplab_G$ is defined as before, and the behavior of the strands at an interior leaf is shown in \cref{fig:interior_leaves}. 
 %
%
\begin{conjecture}\label{conj:top_a_deg}
Let $G$ be a plabic graph with $\conn(G)$ connected components and $n$ interior faces. Then the top $a$-degree of $P(\Lplab_G;a,z)$ equals
 \begin{equation}\label{eq:deg_a}
  %
  \degtop_a(P(\Lplab_G))=\conn(G)-n-1.
 \end{equation}
The degree of $\Ptop_L(q)$ (as a rational function in $q$) is given by
 \begin{equation}\label{eq:a_top}
\deg(\Ptop_L(q)) = n,
 \end{equation}
 and the leading coefficient of $\Ptop_L(q)$ is equal to $1$.
\end{conjecture}

Cluster varieties $\XQice$ are irreducible, with dimension equal to the number of vertices in $\Qice$.  Thus, when $R(Q;q)$ is a rational function, it has leading coefficient $1$ and degree equal to the number of vertices of $Q$.  So \cref{conj:main} implies \eqref{eq:a_top} for simple plabic graphs. %

%
%
%
%
%
%
%
%
%
%
%



%
%
%
%

\subsection{Empty quivers}
\begin{proposition}\label{prop:empty}
  Let $G$ be a connected plabic graph. The following are equivalent:
  \begin{enumerate}
  \item $\QG$ is empty.
  \item $G$ is a tree.
  \item\label{item:empty3} The tail reduction of $G$ is a graph with no interior vertices.
  \end{enumerate}
\end{proposition}
\begin{proof}
The mutable quiver $\QG$ is empty if and only if $G$ has no interior faces, which is equivalent to $G$ being a tree, since $G$ is assumed to be connected.  By Proposition~\ref{prop:trivred}, this is also equivalent to~\eqref{item:empty3}.
\end{proof}

\subsection{Disconnected graphs}
For two plabic graphs $G',G''$, we let $G'\disjun G''$ denote the disconnected plabic graph obtained by taking the disjoint union of $G'$ and $G''$. We denote by the same symbol the disjoint union of two quivers and of two links.
%
\begin{proposition}\label{prop:disc}
Suppose $G=G'\disjun G''$ is a disconnected plabic graph. Then $\QG=\QX(G')\disjun\QX(G'')$ and $\Lplab_G=\Lplab_{G'}\disjun \Lplab_{G''}$.
\end{proposition}
\begin{proof}
Follows directly from the definitions.
\end{proof}

\subsection{Isolated plabic graph quivers}
\newcommand\LHopf{L_{{\rm Hopf}}}
Recall that the \emph{connected sum} of two oriented knots $K, K'$ is defined as follows.  Find planar projections of the two knots that are disjoint.  Then find a (topological) rectangle $R$ in the plane, with oriented boundary consisting of sides $S_1,S_2,S_3,S_4$ in order, such that $S_1$ (resp. $S_3$) is an oriented arc along $K$ (resp. $K'$), and $S_2,S_4$ are disjoint from $K$ and $K'$.  The oriented connected sum knot $K \#K'$ is obtained by deleting $S_1,S_3$ and adding $S_2,S_4$.  To construct the connected sum of two links $L, L'$, we choose components $K \subset L$ and $K' \subset L'$ and perform the above surgery.  We call any link produced in this way the \emph{connected sum} of $L$ and $L'$, denoted $L \#L'$. The following result is well known; see e.g.~\cite[Proposition~16.2]{Lickorish}.

\begin{proposition}\label{prop:FLYsum}
Let $L$ and $L'$ be two oriented links.  Then $\HOMP(L \disjun L') = \left(\frac{a-a^{-1}}{z}\right)\HOMP(L)\HOMP(L')$ and $\HOMP(L \#L') = \HOMP(L) \HOMP(L')$. 
\end{proposition}
In particular, the $r$-component unlink has HOMFLY polynomial $\left(\frac{a-a^{-1}}{z}\right)^{r-1}$.

Let $\LHopf:=\hopflink$ denote the positively oriented Hopf link.  Then we have 
\begin{equation}\label{eq:Hopf}
\HOMP({\LHopf};a,z) = \frac{z + z^{-1}}{a} - \frac{z^{-1}}{a^{3}}.
\end{equation}



Recall that a quiver $Q$ is \emph{isolated} if it has no arrows.
\begin{proposition}\label{prop:plabic_isolated}
Let $G$ be a simple plabic graph with $n$ interior faces and $\conn(G)$ connected components. Suppose that $\QG$ is isolated. Then $\Lplab_G$ is a disjoint union of $\conn(G)$ links, each of which is a connected sum of some number (possibly zero) of positively oriented Hopf links. 
\end{proposition}
\begin{proof}
By Proposition~\ref{prop:disc}, we may assume that $G$ is connected.  Replace $G$ by its tail reduction (Proposition~\ref{prop:trivred}).
Then the result holds when $n = 0$: the link $\Lplab_G$ is an unknot.  The result also holds when $n =1$: the link $\Lplab_G$ is the Hopf link.  So suppose that $n > 1$.  In particular, we are assuming that $G$ has no boundary vertices.

\begin{figure}
\includegraphics[width=1.0\textwidth]{figures/loop_removal}
  \caption{\label{fig:loop_removal} A loop edge in $G$ whose sole vertex is adjacent to a boundary face results in having the Hopf link as a connected summand of $\Lplab_G$; see the proof of \cref{prop:plabic_isolated}.}
\end{figure}

For the remainder of the proof it is convenient to replace $G$ by its \emph{bipartite reduction}, obtained by contracting all interior edges whose endpoints have the same color. In particular, interior vertices are allowed to have degree higher than three.  Suppose first that $G$ contains an interior face $F$ which is bounded by a loop $e$ whose sole endpoint is $v\in V(G)$.  We show that $v$ must be adjacent to a boundary face of $G$. Without loss of generality, we may assume that the edge $e$ is not fully contained inside any other loop edge incident to $v$. Let $F'$ be the face on the other side of $e$ to $F$.  If $F'$ is a boundary face then we are done.  Otherwise, we assume that $F'$ is an interior face.  Let $e'$ be any non-loop edge incident to $v$ that is on the boundary of $F'$.  By our bipartite assumption, $e'$ connects $v$ to a vertex of opposite color.  The assumptions that $G$ is simple and $\QG$ is isolated then imply that $e'$ separates $F'$ from a boundary face.  It follows that $v$ is adjacent to a boundary face of $G$.  Letting $L':=\Lplab_{G-e}$, we see from \cref{fig:loop_removal} that $G$ is a connected sum of $L'$ and $\LHopf$. 
%
  Since $G-e$ has fewer interior faces, the result follows by induction.%

Suppose now that $G$, still assumed to be bipartite, has no interior faces that are bounded by loop edges. Let $F$ be an interior face of $G$. Since $G$ is simple, the perimeter of $F$ does not contain both sides of any edge. Moreover, since $\QG$ is isolated, every edge adjacent to $F$ is adjacent to a boundary face on the other side. Thus, $G$ consists of a number of interior faces connected in a tree-like manner as in \cref{fig:isolated_tree}.  The interior faces at the leaves of this tree must be bounded by loop edges, contradicting our assumption.
  %
\end{proof}

\begin{figure}
  \includegraphics[width=0.4\textwidth]{figures/isolatedtree}
  \caption{\label{fig:isolated_tree} A ``tree'' of interior faces (without boundary vertices) must have a loop.}
\end{figure}

\subsection{Torsion}
%
\newcommand\tF{\tilde F}
By definition, for $G$ a (possibly not simple) plabic graph, the quiver $\Qred_G$ is obtained from the directed graph $\QG$ (\cref{sec:intro:quivers-point-count}) by removing all $1$-cycles, and canceling out directed $2$-cycles one by one.  This agrees with the procedure in \cite{FPST}, and the mutation class of $\Qred_G$ is again invariant under the moves in \cref{fig:moves}. Note that $G$ is simple if and only if $\QG=\Qred_G$.

\begin{proposition}\label{prop:torsionfree}
Let $G$ be a plabic graph. Then $\QredG$ is torsion-free.
\end{proposition}
\begin{proof}
We prove the result by induction on the number of vertices of $\QredG$.  In the following, we index rows of the $\BQredG$ by interior faces $F$ of $G$.

We may suppose that $G$ is connected.  If $G$ has fewer than two interior faces, then checking the graphs in \cref{fig:trivred}, we see that the result holds.  So suppose that $G$ has at least two interior faces.  Applying tail removal, we may suppose that $G$ is not reduced.  Then by \cite{Pos}, one of the following holds: (a) $G$ has an interior leaf, or (b) the bipartite reduction (see the proof of \cref{prop:plabic_isolated}) has a loop, or (c) $G$ is move equivalent to a plabic graph with a two parallel edges between vertices of opposite color.

In case (a), removing the leaf vertex does not change $\QredG$.  We may thus assume that all interior leaves have been removed.  In case (b), deleting the loop edge removes a zero row and a zero column from $\BQredG$, so the result holds by induction.   We may thus assume that $G$ is leafless and loopless and we are in situation (c), that is, $G$ has a double edge: a pair $e',e''$ of edges between vertices $u,v$ of opposite colors.  Let $F$ be the interior face between $e'$ and $e''$, and let $F', F''$ be the faces on the other side of $e'$ and $e''$ to $F$; see \cref{fig:leafface}.   Let $G'$ be obtained from $G$ by removing both edges $e',e''$.  Let $B := \BQredG$, $C := B(\Qred_{G'})$, and $n := |V(\QredG)|$.  If $F',F''$ are both boundary faces, then $\Qred_{G'}$ is obtained from $\Qred_G$ by removing an isolated vertex, so the result holds by induction.  

Suppose that one of the faces $F',F''$ is boundary and one, say $F'$, is interior.    Then $\Qred_{G'}$ is obtained from $\Qred_G$ by removing the vertices $F$ and $F'$.  Let $\Z^V \cong \Z^n$ denote the free $\Z$-module with basis vectors $\{e_\E\}_{\E\in V}$ indexed by the vertex set $V = V(\Qred_G) = \{\mbox{interior faces $\E$ of $G$}\}$ and $\Z^{n-2} \cong \Z^{V'} \subset \Z^V$ the submodule spanned by basis elements other than $e_F, e_{F'}$.  The rows of $B$ belong to $\Z^V$ and the rows of $C$ belong to $\Z^{V'} \subset \Z^V$.  By induction, $C$ is torsion-free, so $\Z^{V'}/C^T \Z^{V'}$ has a basis given by some vectors $d_1,d_2,\ldots,d_r \in \Z^{V'}$.   Furthermore, we have
$$
b_F = \pm e_{F'}, \qquad b_{F'} \pm e_F \in \Z^{V'}, \qquad b_\E - c_\E \in  \Z e_{F'},  
$$
where $\E$ is any interior face other than $F,F'$, and $b_\E$ and $c_\E$ denote the rows of $B$ and $C$, respectively.  It is clear that $\Z^V/B^T \Z^V$ has a basis given by $\{d_1,\ldots,d_r\}$.  So $B$ is torsion-free.

Now suppose that $F',F''$ are both interior faces.  If $F' = F''$, then one of $u,v$, say $v$, must be a cut vertex of $G$.  Thus, $G$ contains a non-trivial induced subgraph $H$ containing $v$ and connected to $G-H$ only via $e',e''$.  The subgraph $H$ is contained ``inside" $F'=F''$; we may consider $H$ to be a plabic graph, necessarily not reduced, with a single boundary vertex incident to $v$.  Replacing $G$ by $H$ and repeating the argument, we are reduced the situation where $F' \neq F''$.

With $F' \neq F''$ both interior, $\Qred_{G'}$ is obtained from $\Qred_G$ by removing the vertex $F$ and identifying the vertices $F'$ and $F''$ to give a new vertex $\tF$ of $\Qred_{G'}$.  The rows of $B$ belong to $\Z^V \cong \Z^n$ and the rows of $C$ belong to $\Z^{V'} \cong \Z^{n-2}$, where $V' = V(\Qred_{G'})$.  Define the inclusion $\iota: \Z^{V'} \hookrightarrow \Z^V$ by sending $e_{\tF}$ to $e_{F'}$, and mapping the other basis vectors in the obvious way.  By induction, $C$ is torsion-free, so $\Z^{V'}/C^T\Z^{V'}$ has a basis given by some vectors $d_1,d_2,\ldots,d_r \in \Z^{V'}$.  We have that 
$$
b_F = \pm(e_{F'} - e_{F''}), \qquad b_{F'} +b_{F''} = \iota(c_{\tF}), \qquad b_\E - \iota(c_\E) \in \Z(e_{F'} - e_{F''}),
$$
where $\E$ is an interior face other than $F,F',F''$.  Let $W \subset \Z^{V - \{F\}}$ be the span of the vectors $b_F, b_{F'} +b_{F''}$ and $b_\E, \E \notin \{F,F',F''\}$.  Then $\{\iota(d_1),\ldots,\iota(d_r)\}$ form a basis of $\Z^{V- \{F\}}/W$.
But $b_{F'} \pm e_F \in\Z^{V - \{F\}}$, so $\Z^V/B^T\Z^V$ has basis $\{\iota(d_1),\ldots,\iota(d_r)\}$.  Thus, $B$ is torsion-free.
\end{proof}
Combining \cref{prop:torsionfree} with \cref{lem:ext}, we obtain the following result.
\begin{corollary}\label{cor:ext_simple}
Let $G$ be a simple plabic graph. Then $\QG$ admits a minimal extension.
\end{corollary}

\section{Skein relations}
In this section, we introduce leaf recurrent plabic graphs.  This class of plabic graphs includes the well-studied reduced plabic graphs and plabic fences.  We show that the leaf recurrence for plabic graphs corresponds to the skein relation for the associated links, thus establishing \cref{conj:main} for leaf recurrent plabic graphs.

\subsection{Leaf recurrent plabic graphs}\label{sec:leaf-recurr-plab}
In this section, we assume that all plabic graphs satisfy the assumptions of Section~\ref{sec:plabquiver}.

\begin{definition}
Let $G$ be a plabic graph. An interior face $F$ of $G$ is a \emph{boundary leaf face} if, in the tail reduction of $G$, the face $F$ has the form shown in \cref{fig:leafface}, i.e.  $F$ is bounded by two edges $e,e'$ with endpoints $u,v$ of distinct color, with $e$ separating $F$ from a boundary face and $e'$ separating $F$ from an interior face $F'$.
\end{definition}

\begin{figure}
  \includegraphics[width=0.7\textwidth]{figures/leafface}
  \caption{\label{fig:leafface} A boundary leaf face $F$ (left) and an interior double edge (right).}
\end{figure}

\begin{lemma}\label{lemma:bdry_leaf_face}
Let $G$ be a simple plabic graph and $F$ a boundary leaf face. Let $u,v\in V(G)$ be the vertices adjacent to $F$. Then the plabic graphs $G':=G-e$ and $G'':=G-\{u,v\}$ are both simple.
\end{lemma}
\begin{proof}
The directed graph $\QX(G')$ is obtained from $\QG$ by removing the vertex corresponding to $F$.  Similarly, $\QX(G'')$ is obtained from $\QG$ by removing the vertices corresponding to $F$ and $F'$.  Both directed graphs are clearly quivers, and thus $G',G''$ are simple.
\end{proof}

\begin{definition}\label{dfn:leaf_recurrent_plabic_graph}
  The class of \emph{leaf recurrent} plabic graphs is defined as follows.
  \begin{enumerate}[(a)]%
  \item Every leaf recurrent plabic graph is simple.
  \item If $\QG$ is isolated then $G$ is leaf recurrent.
  \item If $G$ is move equivalent to a leaf recurrent plabic graph then $G$ is leaf recurrent.
  \item Suppose $G$ has a boundary leaf face $F$ as in \cref{lemma:bdry_leaf_face}. If the plabic graphs $G'$ and $G''$ are leaf recurrent then $G$ is leaf recurrent.
  \end{enumerate}
\end{definition}

\begin{remark}\label{rmk:leaf_rec=>Louise}
It is immediate that if $G$ is leaf recurrent then $\QG$ is \Louise.
\end{remark}

\begin{theorem}[{\cite[Remark~4.7]{MSLA}}]\label{thm:MSLA}
Any reduced plabic graph $G$ is leaf recurrent.
\end{theorem}



\subsection{The skein relation}\label{sec:sub_skein-relation}
%



\begin{figure}
  \includegraphics[width=1.0\textwidth]{figures/skein}
  \caption{\label{fig:skein} Leaf recurrence for plabic graphs corresponds to the \FLY skein relation; see \cref{lemma:bdry_leaf_skein_FLY}.}
\end{figure}

\begin{lemma}\label{lemma:bdry_leaf_skein_FLY}
Suppose that $G$ is a plabic graph with a boundary leaf face $F$. Let $L_+:=\Lplab_G$, $L_0:=\Lplab_{G'}$, and $L_-:=\Lplab_{G''}$, where $G,G',G''$ are the three graphs in \cref{lemma:bdry_leaf_face}. Then the \FLY polynomials of $(L_+,L_0,L_-)$ are related by~\eqref{eq:HOMFLY_dfn}.
\end{lemma}
\begin{proof}
See \cref{fig:skein}.
\end{proof}

\begin{theorem}
  Suppose $G$ is a leaf recurrent plabic graph. Then Conjectures~\ref{conj:main} and~\ref{conj:top_a_deg} hold for $G$.
\end{theorem}
\begin{proof}
  Suppose $G$ has a boundary leaf face $F$ such that the graphs $G'$, $G''$ from \cref{lemma:bdry_leaf_face} are both leaf recurrent. Denote $Q:=\QG$, $Q':=\QX(G')$, $Q'':=\QX(G'')$, $L_+:=\Lplab_G$, $L_0:=\Lplab_{G'}$, and $L_-:=\Lplab_{G''}$, and note that $G,G',G''$ all have the same number of connected components.

By \cref{lemma:bdry_leaf_skein_FLY}, the \FLY polynomials of $(L_+,L_0,L_-)$ are related by~\eqref{eq:HOMFLY_dfn}. By \cref{rmk:leaf_rec=>Louise}, the quivers $Q,Q',Q''$ are \Louise. By \cref{cor:RQ_skein}, their point counts are related by~\eqref{eq:RQ_skein}. Comparing~\eqref{eq:HOMFLY_dfn} and~\eqref{eq:RQ_skein}, we see that~\eqref{eq:FLY=pcnt} holds for $Q$ if it holds for both $Q'$ and $Q''$.

Now suppose that $\QG$ is isolated with $n$ vertices.  If $G$ is the disjoint union of $G_1$ and $G_2$ then by Proposition~\ref{prop:FLYsum}, we have $\HOMP(\Lplab_{G})= ((a-a^{-1})/z)\HOMP(\Lplab_{G_1})\HOMP(\Lplab_{G_2})$, so $\Ptop_{\Lplab_G}(q) = (q-1)^{-1} \Ptop_{\Lplab_{G_1}}(q)\Ptop_{\Lplab_{G_2}}(q)$.  Since $\RQG(q) = \RQX{\QX(G_1)}{q}\RQX{\QX(G_2)}{q}$, the result follows. We may now assume that $G$ is connected.  By \cref{prop:plabic_isolated}, $\Lplab_G$ must be the connected sum of $n$ Hopf links.  By \cref{prop:FLYsum}, \eqref{eq:Hopf} and~\eqref{eq:A1} we have 
$$\Ptop_{\Lplab_G}(q) = \left(\frac{q^2-q+1}{q-1}\right)^n = \RQG(q).
$$
Thus, the result holds when $\QG$ is isolated.

Since the validity of~\eqref{eq:FLY=pcnt} is unchanged under moves on $G$, it follows that~\eqref{eq:FLY=pcnt} holds for all leaf recurrent plabic graphs $G$.  The equality~\eqref{eq:deg_a} also follows from the same argument.
\end{proof}

\subsection{Plabic fences}\label{sec:plabic-fences}
Following~\cite[Section~12]{FPST}, we consider the following class of plabic graphs. Let $\n\geq 2$ and let $I:=[\n-1]=\{1,2,\dots,\n-1\}$. A \emph{double braid word} is a word $\bw=(i_1,i_2,\dots,i_m)$ in the alphabet
\begin{equation*}%
  \pmn:=\{-1,-2,\dots,-(\n-1)\}\sqcup\{1,2,\dots,\n-1\}.
\end{equation*}
We denote the set of double braid words of length $m$ by $\DRW$. 

To a double braid word $\bw\in\DRW$ we associate a plabic graph $\Gbw$ with $2\n$ boundary vertices. First, draw $\n$ horizontal strands, with both endpoints of each strand on the boundary of the disk. Then for each $j=1,2,\dots,m$, let $h:=|i_j|$ and consider the strands at height $h$ and $h+1$. If $i_j>0$, add a black at $h$, white at $h+1$ bridge between these two strands, and if $i_j<0$, add a white at $h$, black at $h+1$ bridge between these two strands. The bridges are added one by one proceeding from left to right. See \cref{fig:PF2} for an example.

\begin{figure}
  \includegraphics[width=0.7\textwidth]{figures/plabic_fence-2}
  \caption{\label{fig:PF2} The plabic fence $\Gbw$ for a double braid word $\w$.}
\end{figure}

\begin{proposition}
For any $\bw\in\DRW$, the plabic graph $\Gbw$ is simple.
\end{proposition}
\begin{proof}
This is clear from construction. %
\end{proof}

Let us say that two braid words $\bw,\bw'$ are \emph{double braid move equivalent} if they are related by a sequence of the following moves:
\begin{enumerate}[(B1)]%
\item\label{bm1} $(\dots,i,j,i,\dots)\leftrightarrow (\dots,j,i,j,\dots)$ \quad if $i,j\in\pmn$ have the same sign and $|i-j|=1$;
\item\label{bm2} $(\dots,i,j,\dots)\leftrightarrow (\dots,j,i,\dots)$ \quad if $i,j\in \pmn$ have the same sign and $|i-j|>1$;
%
%
\item\label{bm5} $(\dots,i,j,\dots)\leftrightarrow (\dots,j,i,\dots)$  \quad if $i, j\in\pmn$ have different signs;
\item\label{bm6} $(-i,\dots)\leftrightarrow (i,\dots)$ \quad for $i\in \pmn$;
\item\label{bm7} $(\dots,-i)\leftrightarrow (\dots,i)$ \quad for $i\in \pmn$.
\end{enumerate}

%
%
%
%
%
It is clear that if $\bw,\bw'\in\DRW$ are double braid move equivalent then the plabic graphs $\Gbw,\GX(\bw')$ are  move equivalent.

\begin{proposition}
Plabic fences are leaf recurrent. That is, for any double braid word $\bw\in\DRW$, $\Gbw$ is a leaf recurrent plabic graph.
\end{proposition}
\begin{proof}
Proceed by induction on $m$. The base case $m=0$ is clear.

Applying moves~\ref{bm5}--\ref{bm6}, we may assume that $\bw$ has no negative indices. We say that a word $\bw=(i_1,i_2,\dots,i_m)\in I^m$ is \emph{reduced} if the associated permutation $\pi(\bw):=s_{i_1}s_{i_2}\cdots s_{i_m}$ has precisely $m$ inversions. If $\bw$ is not reduced then we may apply moves~\ref{bm1}--\ref{bm2} to transform $\bw$ into a word of the form $(\bw_1,i,i,\bw_2)$ for some $i\in I$ and some words $\bw_1\in I^{m_1}$, $\bw_2\in I^{m_2}$. Let $-\rev(\bw_1)\in (-I)^{m_1}$ denote the word obtained from $\bw_1$ by negating all indices and reversing their order. One can transform $\bw_1$ into $-\rev(\bw_1)$ by applying moves~\ref{bm5}--\ref{bm6}. Thus, applying moves~\ref{bm5}--\ref{bm6}, we transform
\begin{equation*}%
  (\bw_1,i,i,\bw_2)\to (-\rev(\bw_1),i,i,\bw_2)\to (i,i,\bw_2,-\rev(\bw_1)) \to (i,i,\bw_2,\bw_1).
\end{equation*}  The double bridge in $\GX(i,i,\bw_2,\bw_1)$ corresponding to $(i,i,\dots)$ gives rise to a boundary leaf face. Applying the induction hypothesis to the words $\bw':=(i,\bw_2,\bw_1)$ and $\bw'':=(\bw_2,\bw_1)$, we see that the plabic graphs $\GX(\bw')$ and $\GX(\bw'')$ are leaf recurrent. Thus, by \cref{dfn:leaf_recurrent_plabic_graph}, $\GX(\bw)$ is also a leaf recurrent plabic graph.

Finally, suppose that the word $\bw=(i_1,i_2,\dots,i_m)\in I^m$ is reduced. Then $\Gbw$ is a reduced plabic graph and therefore we are done by \cref{thm:MSLA}. Alternatively, applying moves~\ref{bm5}--\ref{bm6}, we may transform $\bw$ as follows:
\begin{equation*}%
  (i_1,i_2,\dots,i_m) \to (-i_1,i_2,\dots,i_m) \to (i_2,\dots,i_m,-i_1)\to (i_2,\dots,i_m,i_1).
\end{equation*}
We refer to this transformation as a \emph{conjugation move}. Applying conjugation and braid moves repeatedly, we either arrive at a non-reduced word and proceed by induction, or we find that $\pi(\bw)\in\Sn$ is a minimal length representative in its conjugacy class; see~\cite[Theorem~1.1]{GePf}. Recall that conjugacy classes in the symmetric group $\Sn$ correspond to cycle types, and $\pi(\bw)$ has minimal length  if and only if it is a product of cycles of the form $(a,a+1,\dots,b)$ on disjoint intervals $[a,b]\subset[\n]$. The resulting plabic fence $\Gbw$ has no interior faces and therefore is leaf recurrent.
\end{proof}
